\section{Economic rationality: conditional representations and contrasts}
\label{sec:p19-annex-rationality-9-4}

\begin{annexitem}{annex:p19.r94.continuous-converse}{thm:p19-r94-continuous-converse}{Continuous representation and nonattainment}
\begin{proof}
For any two real numbers \(u(x),u(y)\), at least one of
\(u(x)\geq u(y)\) and \(u(y)\geq u(x)\) holds; this proves
completeness. If \(x\succeq y\succeq z\), the representation gives
\(u(x)\geq u(y)\geq u(z)\), hence \(x\succeq z\).
For fixed \(y\), the upper weak contour is the inverse image of
\([u(y),\infty)\) under \(u\), and the lower weak contour the
inverse image of \(( -\infty,u(y)]\). Both are closed in the relative
topology of \(X\) by continuity.

For a strictly increasing function \(g\) on the range of \(u\),
\(u(x)\geq u(y)\) holds exactly when
\(g(u(x))\geq g(u(y))\); thus \(g\circ u\) preserves every
comparison, including equality. No continuity claim about an arbitrary
\(g\) is needed for this ordinal conclusion.
Finally, for every \(x\in(0,1)\), \((x+1)/2\) belongs to the domain
and is strictly larger. The identity utility is continuous but has no
maximizing point there.
\end{proof}
\end{annexitem}

\begin{annexitem}{annex:p19.r94.source-vnm}{thm:p19-r94-source-vnm}{The calibrated finite-prize representation}
\begin{proof}
After reduction, lotteries are probability vectors. Put
\(B_q=qM+(1-q)m\). The replacement and strict segment-ordering
argument of \cref{thm:p17-mixture-comparison} uses only a weak order
and common-mixture independence; it therefore applies here. In particular,
\(B_q\succ B_r\) exactly when \(q>r\), and indifferent positive-
weight component replacements preserve indifference. Zero weights leave
the vector unchanged. No topological continuity hypothesis is used in
these facts.

Endpoint calibration makes each certain prize lottery \(\delta_{x_i}\)
indifferent to some \(B_{u_i}\). Strict segment ordering makes \(u_i\)
unique, with \(0\leq u_i\leq1\), and gives 1 and 0 at the selected
endpoints. For any lottery with probabilities \(p_i\), its reduced
decomposition is \(\sum_i p_i\delta_{x_i}\). Replacing each component
by its indifferent endpoint mixture, one at a time, yields
\[
 \sum_i p_i\delta_{x_i}\sim\sum_i p_i B_{u_i}
   =\left(\sum_i p_i u_i\right)M+
                  \left(1-\sum_i p_i u_i\right)m.
\]
The equality is coordinatewise, using \(\sum_i p_i=1\); reduction
identifies the compound procedures with these vectors. Indifference
replacements and strict segment ordering show that the weak comparison
between any two lotteries holds exactly when their \(\sum_i p_i u_i\)
values have the corresponding inequality. This proves representation and
normalized uniqueness.

If another prize function \(v\) represents the same relation by expected
utility, write \(v(M)\) and \(v(m)\) for its values on the certain
prizes underlying those lotteries. Then \(v(M)>v(m)\). Equality of expectations of
\(\delta_{x_i}\) and \(B_{u_i}\) forces
\[
             v(x_i)=v(m)+(v(M)-v(m))u_i.
\]
Thus it is a positive-affine transform; every such transform preserves
the preference because lottery probabilities sum to one.

Conversely, suppose prize utility represents a weak order by expectation,
with designated best utility strictly larger than worst utility and all
prize utilities between them. Real comparison is complete and transitive.
Each expectation lies between the endpoint values, so solving the linear
endpoint-mixture equation gives calibration in \([0,1]\). Mixing two
compared expectations with a common third multiplies their difference by
\(q>0\), preserving its sign and zeros. Reduction preserves expectation
by distributing finite sums. These verify all stated axioms, including
endpoint and zero-branch cases.
\end{proof}
\end{annexitem}

\begin{annexitem}{annex:p19.r94.backward}{thm:p19-r94-backward}{Finite backward induction and the two-date calculation}
\begin{proof}
At a solvent terminal state after date \(T\), no utility remains, so
value is zero. At the final decision date, the finite nonempty feasible
set has a maximizer of current utility among choices leading to those
terminal states. Suppose values and maximizing continuations have been
constructed for date \(i+1\). Every feasible date-\(i\) plan begins
with a feasible action \(c\), reaches its specified next state, and has
discounted continuation no larger than \(\delta V_{i+1}\) there.
Its value is no larger than that action's displayed maximand, hence no
larger than \(V_i\). Conversely, a maximizing action exists in the finite
set. Follow it by a continuation attaining its next-state value, supplied
by the induction hypothesis, to attain \(V_i\). Backward induction
through finitely many dates proves value and attainment. Insolvent or
continuationless branches are removed before maximizing, so no empty
maximum or infeasible policy is used.

Multiplying the date-\(i+1\) continuation by \(\delta\) changes its
exponents to \(j-i\), the forward discounting convention in the body.
In the example, increasing utility makes final exhaustion optimal: any
leftover integer stock can be consumed at date 2 to increase utility.
Exhaustion leaves exactly the three displayed allocations.
Since \(\sqrt2<3/2\) (square both positive sides to obtain
\(2<9/4\)) and \(\sqrt2/2<\sqrt2\), \(3/2\) is largest.
\end{proof}
\end{annexitem}

\begin{annexitem}{annex:p19.r94.time-contrast}{thm:p19-r94-time-contrast}{Common shifts and a present-bias reversal}
\begin{proof}
Adding the same admissible date shift \(h\) to exponential scores
multiplies both by \(\delta^h>0\), preserving the sign of their
difference, including zero.

For linear utility and zero reward at other dates, 10 now scores 10,
while 15 one date later scores \((1/2)(9/10)15=27/4=6.75\).
Thus immediate 10 is preferred. Viewed from date 0, rewards 10 at date 1
and 15 at date 2 score \((1/2)(9/10)10=9/2=4.5\) and
\((1/2)(9/10)^2 15=243/40=6.075\), respectively, favoring the latter.
Viewed from date 1, they are 10 now and 15 one date later, so the first
calculation gives the opposite ranking. Only the evaluation date changes.
This proves the common-shift and re-evaluation contrast for the declared
model, without any population-frequency assertion.
\end{proof}
\end{annexitem}

\begin{annexitem}{annex:p19.r94.reversal}{thm:p19-r94-reversal}{The precise preference-reversal contradiction}
\begin{proof}
Indifference gives \(\$\text{-bet}\succeq C_{\$}\) and
\(C_P\succeq P\text{-bet}\). The comparison
\(C_{\$}\succ C_P\) includes \(C_{\$}\succeq C_P\).
Transitivity therefore gives \(\$\text{-bet}\succeq P\text{-bet}\).
But \(P\text{-bet}\succ\$\text{-bet}\) explicitly denies that
weak comparison, a contradiction. Under strict monetary monotonicity with
other consequences fixed, \(C_P<C_{\$}\) supplies the required strict
comparison. Without that premise numeric inequality alone does not do so.
\end{proof}
\end{annexitem}

\begin{annexitem}{annex:p19.r94.allais}{thm:p19-r94-allais}{Cancellation in the declared Allais contrast}
\begin{proof}
Write \(u_0,u_1,u_5\) for utility of the three prizes. The first strict
comparison requires
\[
    u_1>.89u_1+.10u_5+.01u_0,
       \quad\text{hence}\quad .11u_1>.10u_5+.01u_0.
\]
The second requires
\[
    .10u_5+.90u_0>.11u_1+.89u_0,
       \quad\text{hence}\quad .10u_5+.01u_0>.11u_1.
\]
These are opposite strict inequalities between the same real numbers.
No \(u\) satisfies both. All lottery probabilities are nonnegative
and sum to one. The contradiction uses linearity in probabilities, not
linearity of prize utility.
\end{proof}
\end{annexitem}
